\begin{conclusion}

	本文建立了一个 SEIRV 仓室模型，该模型考虑了年龄结构和疫苗接种，并应用于分析2021年初石家庄市 COVID-19 疫情的流行病学和动力学特征。

	研究首先关注了疫情的流行病学特征，包括时间序列、年龄-性别分布和地理分布。这些因素被纳入分析，以深入了解疫情的传播模式和影响。
	其次，对微分方程模型的动力学性质进行了研究。模型的性质包括无病平衡点的存在唯一性定理以及这些平衡点的稳定性定理。此外，还探讨了地方病平衡点的存在唯一性定理和稳定性定理。这些理论结果有助于理解疫情传播的基本特征。
	最后，对神经网络模型的动力学特征进行了考察，包括控制措施开始前的失觉期、疫情中的拐点以及最终规模的预测。这些特征有助于制定有效的防控策略和预测疫情的发展趋势。
	通过 SEIRV 仓室模型和神经网络模型的结合，以及对流行病学和动力学特征的深入分析，为更好地理解和应对 COVID-19 疫情提供了重要的理论基础和实际支持。
\end{conclusion}